\section{Budget--Stability Probe: Procedure and Superseded Findings}
\label{app:bsweep}

% =========================================================================
% This appendix accompanied the former pilot. Like the pilot section, it has
% been neutralized: the budget--stability sweep ran against the superseded
% action space and search configuration of the early implementation, so its
% per-cell numbers no longer describe the system this paper specifies. We
% retain the procedure and the qualitative lessons that survive the redesign,
% because they justify pre-registered choices in Section~\ref{sec:experimental-plan};
% the numeric table and curves are not carried forward.
% =========================================================================

This appendix records the procedure of the budget--stability probe summarized in
Section~\ref{sec:pilot-bsweep} and the qualitative lessons that survive the
action-space redesign of Section~\ref{sec:method-actions}. It reports no per-cell
numbers: the probe ran against the superseded verb set and a search
configuration below the pre-registered budgets, so its specific agreement,
sharpness, and distance values describe a system this paper no longer specifies.
What carries forward is the design of the diagnostic and the way it informs the
pre-registered plan.

\paragraph{Per-call seeding does not stabilize a collection.} The stochastic LLM
transitions in the probe were drawn at a nonzero temperature with no fixed
sampling seed. A natural first attempt at reproducibility is to forward a
per-call seed (derived from the planner's random state) to the inference server,
which both Ollama and vLLM accept. We implemented this and verified that an
\emph{individual} call reproduces at a matched seed. It does not, however, make a
\emph{collection} reproducible, and the reason is structural rather than
numeric: an MCTS search's sequence of LLM calls is itself a function of the
responses, because each response is the next state and the next state determines
the next prompt. A single non-deterministic call early in the search
(server-side batching and floating-point reduction order make exact determinism
fragile under load) changes which node expands next, desynchronizes the
seed-to-prompt correspondence, and every downstream call then lands on a
different prompt. One can seed the calls but not the tree they induce; the
seeding hook is retained because it makes single transitions reproducible for
debugging, not because it stabilizes collection. This is a structural finding
about the search-induced prompt cascade and survives the redesign unchanged.

\paragraph{Sweep design.} For a fixed initial state we ran many independent MCTS
searches (different random seeds, hence independent stochastic LLM transitions)
at a range of simulation budgets, on a low-ambiguity control problem and a
high-ambiguity problem, with the grid deliberately spanning the low budget the
probe collected with up through substantially larger budgets. Holding the start
state fixed across all runs isolates the visit-count estimator's own variance
from the trajectory cascade: any disagreement among the resulting $\pimcts$
distributions is the estimator's sampling noise at that budget, not a divergent
trajectory. The stability statistic is \emph{pairwise argmax agreement}: the
probability that two independent searches recommend the same verb, computed over
all pairs of runs, with bootstrap percentile confidence intervals. The
per-budget numeric values from the early sweep are superseded and are not
reproduced here.

\paragraph{Qualitative findings that survive.} Four lessons carry into the
pre-registered study; none depends on the superseded numbers.

First, \emph{sharpness rises with the search budget}: as the budget grows the
visit-count estimator concentrates on the high-value verb and the distribution's
entropy falls, exactly as a sampled-Bellman argument predicts.

Second, \emph{run-to-run agreement on the recommended verb rises with the
budget, and the low budget the probe collected with sits near the agreement
floor}. Convergence speed is family-dependent: the low-ambiguity family
stabilizes at a smaller budget than the high-ambiguity one. The operational
consequence is that the heavy-search budget must be set by the
slowest-converging family rather than the average, which justifies the
pre-registered $N_{\mathrm{large}} = 500$ sitting well above the probe's budget.

Third, \emph{mean $L_1$ between visit-count distributions is a confounded
stability metric and pairwise recommended-verb agreement is the correct one}. As
the budget grows the distributions sharpen, so a single sharp minority run on a
different verb inflates $L_1$ even as true run-to-run agreement improves; on the
high-ambiguity family $L_1$ can rise with the budget while agreement also rises,
so the metric punishes exactly the sharpening that stability requires. The
pre-registered target-stability diagnostic is therefore pairwise
recommended-verb agreement across independent searches from a fixed state,
reported stratified by problem family (Section~\ref{sec:pilot-lessons}).
Relatedly, the recommended verb can itself be budget-indexed on the hard family:
deeper trees back up different values, so the modal verb can shift between
budgets even as within-budget agreement rises. ``Stable target'' therefore means
stable \emph{at a fixed budget}, and any round-over-round comparison must hold
the budget constant.

Fourth, \emph{small-sample stability diagnostics can manufacture qualitative
structure}. An early sweep with few independent searches per cell showed an
apparent non-monotonicity in the budget--stability curve that dissolved once the
number of searches per cell was raised. Conclusions about the shape of the curve
need a generous number of independent searches per cell, well above what a quick
probe suggests.

\paragraph{Cost and the carried-forward commitments.} Heavy search is expensive:
a single large-budget search took on the order of minutes on the probe's
hardware, so the pre-registered $N_{\mathrm{large}} = 500$ implies a
correspondingly costly search per decision. The convergence trend is the
justification for paying it. Two practical commitments follow for the larger
study, beyond those in Section~\ref{sec:pilot-lessons}: report the
pairwise-agreement diagnostic \emph{stratified by problem family} rather than
pooled, and set the heavy-search budget by the slowest-converging family rather
than the average.
